% Lösungen Potenzen, Zehnerpotenzen und Quadratwurzeln · Quadratwurzeln · Prüfungs-Fokus MSA (zusammenbau v0.6)
\einheitenkopf{Einheit 3 · Quadratwurzeln}

\erg{1}{a) ja, Quadratzahl – elf mal elf \quad b) $7$, Probe: $7 \cdot 7 = 49$ \quad c) $(-8)^2 = 64$, $\sqrt{64} = 8$}
\erg{2}{a) $\sqrt{6{,}25} = 2{,}5$, die Seite ist $2{,}5$ m lang \quad b) $\sqrt{1\,600} = 40$, die Seite ist $40$ m lang \quad c) $\sqrt{441} = 21$, die Seite ist $21$ cm lang}
\erg{3}{a) $\sqrt{(-7)^2} = 7$, denn erst das Quadrat: $(-7)^2 = 49$, und die Wurzel ist nie negativ \quad b) $\sqrt{(-13)^2} = 13$, denn erst das Quadrat: $(-13)^2 = 169$, und die Wurzel ist nie negativ \quad c) $\frac{7}{3} > \frac{9}{4}$, denn $\frac{7}{3} \approx 2{,}33$ und $\frac{9}{4} = 2{,}25$; $\sqrt{5} \approx 2{,}24$ ist kleiner \quad d) $\sqrt{3} < \frac{7}{4}$, denn $\sqrt{3} \approx 1{,}73$ und $\frac{7}{4} = 1{,}75$ \quad e) $-0{,}26 < -\frac{1}{4} < 1{,}7 < \sqrt{3}$, denn $-\frac{1}{4} = -0{,}25$ und $\sqrt{3} \approx 1{,}73$ \quad f) $-0{,}62 < -\frac{3}{5} < 3{,}3 < \sqrt{11}$, denn $-\frac{3}{5} = -0{,}6$ und $\sqrt{11} \approx 3{,}32$ \quad g) $\sqrt{7{,}5^2 - 2{,}1^2} = \sqrt{56{,}25 - 4{,}41} = \sqrt{51{,}84} = 7{,}2$, die Leiter reicht $7{,}2$ m hoch \quad h) $\sqrt{4^2 - 3{,}6^2} = \sqrt{16 - 12{,}96} = \sqrt{3{,}04} \approx 1{,}74$, der Höhenunterschied ist etwa $1{,}74$ m \quad i) $r = \sqrt{750 : (\pi \cdot 9)} = \sqrt{26{,}53\ldots} \approx 5{,}15$, der Radius ist etwa $5{,}15$ cm \quad j) $330$ ml $= 330\,\text{cm}^3$; $r = \sqrt{330 : (\pi \cdot 8)} = \sqrt{13{,}13\ldots} \approx 3{,}62$, der Radius ist etwa $3{,}62$ cm \quad k) $x = 4 \pm \sqrt{16 - 13} = 4 \pm \sqrt{3}$, also $x_1 \approx 5{,}73$ und $x_2 \approx 2{,}27$ \quad l) $x = -2 \pm \sqrt{4 + 1} = -2 \pm \sqrt{5}$, also $x_1 \approx 0{,}24$ und $x_2 \approx -4{,}24$}
